\documentclass[12pt]{report}
\usepackage[T1]{fontenc}
\usepackage[brazilian]{babel}
\usepackage{csquotes}
\usepackage{makeidx}
\usepackage{glossaries}
\usepackage[
  perfil=academico,
  tipo-trabalho=monografia,
  impressao=frente-e-verso,
  citacoes=autor-data
]{abntex3}

\makeindex
\makeglossaries
\newglossaryentry{metadado}{
  name={metadado},
  description={dado estruturado usado para descrever outro recurso}
}
\addbibresource{referencias-exemplo.bib}
\ABNTEXsetup{
  titulo={Modelo canônico completo de monografia},
  subtitulo={base auditável com todos os elementos implementados},
  autoria={Maria da Silva},
  instituicao={Universidade Exemplo},
  natureza={Monografia de especialização},
  objetivo={apresentada para obtenção do título de Especialista},
  area={Gestão da Informação},
  orientacao={Prof. Dr. João Souza},
  coorientacao={Profa. Dra. Ana Santos},
  local={Natal},
  data={2026}
}
\ABNTEXabreviatura{p.}{página}
\ABNTEXsigla{PDF}{Portable Document Format}
\ABNTEXsimbolo{$x$}{Valor observado}

\begin{document}
\ABNTEXacademicoexterna
\ABNTEXlombada[Monografia 2026][Universidade Exemplo]

\ABNTEXacademicopretextual
\ABNTEXfolhaderosto
\ABNTEXfichacatalografica{%
  Ponto de extensão para a ficha catalográfica, quando aplicável.
}
\ABNTEXerrata
  {SILVA, Maria da. \textbf{Modelo canônico completo de monografia}.
   2026. Monografia — Universidade Exemplo, Natal, 2026.}
  {\begin{center}
     \begin{tabular}{llll}
       Folha & Linha & Onde se lê & Leia-se\\
       7 & 6 & metadadoso & metadado
     \end{tabular}
   \end{center}}
\ABNTEXfolhadeaprovacao
  {30 de julho de 2026}
  {%
    \ABNTEXmembrobanca
      {Prof. Dr. João Souza}{Orientador}{Universidade Exemplo}
    \ABNTEXmembrobanca
      {Profa. Dra. Ana Santos}{Coorientadora}{Universidade Exemplo}
    \ABNTEXmembrobanca
      {Profa. Ma. Beatriz Alves}{Examinadora}{Instituto Exemplo}
  }
\ABNTEXdedicatoria{Às pessoas que ensinam e aprendem continuamente.}
\ABNTEXagradecimentos{À instituição e às pessoas que avaliaram o trabalho.}
\ABNTEXepigrafe[Autoria demonstrativa]{Clareza favorece a revisão.}
\ABNTEXresumo
  {Modelo completo de monografia que demonstra os elementos públicos,
   alternativas selecionadas e pontos destinados à auditoria.}
  {monografia, metadados, auditoria, LaTeX}
\ABNTEXresumoemlinguaestrangeira
  {Complete monograph template demonstrating public elements, selected
   alternatives, and external audit points.}
  {monograph, metadata, audit, LaTeX}
\ABNTEXlistadeilustracoes
\ABNTEXlistadetabelas
\ABNTEXlistadeabreviaturasesiglas
\ABNTEXlistadesimbolos
\ABNTEXsumario

\ABNTEXacademicotextual
\chapter{Introdução}
\index{metadado}
Apresentam-se contexto, problema, objetivos e justificativa. Um
\gls{metadado} exemplifica o glossário \parencite{associacao2025}.
\section{Questão de estudo}
Como a explicitação da estrutura favorece a avaliação?
\section{Objetivos}
Demonstrar a estrutura e produzir uma referência reutilizável.

\chapter{Revisão de literatura}
\textcite{associacao2025} demonstra a chamada autor-data integrada.
\begin{ABNTEXcitacaolonga}
Trecho criado exclusivamente para verificar o tratamento visual de uma
citação longa no modelo canônico.
\end{ABNTEXcitacaolonga}
\ABNTEXnota{Nota demonstrativa da monografia.}

\chapter{Método}
\section{Procedimentos}
O método combina inventário técnico, compilação e inspeção.

\chapter{Análise e discussão}
\begin{figure}[ht]
  \centering
  \rule{0.55\linewidth}{2cm}
  \caption{Estrutura conceitual demonstrativa}
  \ABNTEXfonte{elaborada pela autora.}
\end{figure}
\begin{table}[ht]
  \centering
  \caption{Itens inspecionados}
  \begin{tabular}{ll}
    Item & Evidência\\
    Elementos & presentes\\
    Ordem & verificável
  \end{tabular}
  \ABNTEXfonte{elaborada pela autora.}
\end{table}

\chapter{Considerações finais}
Retomam-se a questão, os objetivos, as limitações e recomendações.

\ABNTEXacademicopostextual
\ABNTEXbibliografia
\ABNTEXglossario
\ABNTEXappendix{Quadro elaborado pela autora}
Material produzido para apoiar a análise.
\ABNTEXannex{Documento externo de demonstração}
Material não elaborado pela autoria.
\ABNTEXindice[Índice de assuntos]
\end{document}
